\chapter{Fact Table dan Skema Multi-Fakta}
\label{modul:fact-table-multi-fakta}

\section{Identitas Modul}

\begin{identitasmodul}
\textbf{Sumber materi}    & Pertemuan 5 \\
\textbf{CPMK}             & CPMK-092 \\
\textbf{Minggu pelaksanaan} & Minggu ke-6 \\
\textbf{Durasi tatap muka} & 120 menit \\
\textbf{Estimasi tugas}   & 3--4 jam di luar sesi \\
\textbf{Moda kerja}       & Individual \\
\textbf{Perangkat}        & PostgreSQL 17 dan DBeaver \\
\textbf{Dataset}          & Penjualan, persediaan, dan promosi NusaMart \\
\textbf{Skrip pendukung}  & \berkas{sql/modul-04/} beserta \berkas{check.sql} \\
\textbf{Luaran}           & Tiga fact table beserta pernyataan grain dan satu laporan drill-across \\
\end{identitasmodul}

\slideref{5}

\section{Capaian dan Indikator Keberhasilan}

Mahasiswa mampu memilih jenis fact table berdasarkan perilaku proses,
menyatakan grain setiap fakta, serta menyandingkan beberapa fakta melalui
dimensi bersama. Drill-across harus dilakukan setelah agregasi masing-masing
fakta.

\begin{capaian}
Pada akhir modul ini mahasiswa mampu:
\begin{enumerate}[leftmargin=1.4em]
  \item membedakan empat jenis fact table dan memilihnya dari perilaku proses,
        bukan dari bentuk datanya;
  \item membangun \istilah{periodic snapshot} beserta measure semi-aditif dan
        menjelaskan mengapa measure itu tidak boleh dijumlahkan lintas waktu;
  \item membangun \istilah{factless fact} untuk mencatat kejadian yang tidak
        memiliki measure numerik;
  \item menyandingkan dua proses lewat dimensi terkonformasi dengan
        \istilah{drill-across}, bukan dengan men-join dua tabel fakta;
  \item mengenali dan menjelaskan penggandaan measure yang timbul bila aturan
        di atas dilanggar.
\end{enumerate}
\end{capaian}

Sampai Modul 3, gudang data Anda hanya memuat satu proses bisnis. Mulai modul
ini ia memuat tiga --- dan persoalan baru muncul: bagaimana menjawab pertanyaan
yang memerlukan dua proses sekaligus, tanpa membuat angkanya menggelembung.

\section{Prasyarat dan Persiapan}

\subsection{Prasyarat}

\begin{itemize}
  \item Skema bintang hasil Modul 3, termasuk \perintah{dim\_produk} yang sudah
        SCD-2 dan \perintah{dim\_promosi};
  \item pengertian grain dan sifat keteraditifan measure dari Modul 2;
  \item bus matrix hasil tugas Modul 1 --- kolom dimensi yang dipakai bersama
        oleh beberapa proses;
  \item agregasi SQL dengan \perintah{CTE}, serta pengertian
        \perintah{FULL OUTER JOIN} dan \perintah{COALESCE}.
\end{itemize}

\subsection{Pre-lab}

\begin{enumerate}[leftmargin=1.4em]
  \item Bila Modul 3 belum selesai, pulihkan \berkas{snapshot-modul-03.sql}.
  \item Periksa bahwa ketiga dimensi bersama sudah ada dan terisi:
        \perintah{dim\_tanggal}, \perintah{dim\_produk}, dan
        \perintah{dim\_toko}.
  \item Buka kembali bus matrix Modul 1. Tandai dimensi mana yang dipakai oleh
        ketiga proses yang akan dibangun hari ini --- penjualan, persediaan, dan
        promosi.
  \item Siapkan jawaban berikut pada \berkas{modul-04-multi-fakta/pre-lab.md}:
  \begin{enumerate}[label=\alph*.,leftmargin=1.4em]
    \item Saldo persediaan sebuah produk pada 1--3 Januari berturut-turut 100,
          120, dan 90 unit. Berapa ``total saldo'' selama tiga hari itu, dan
          apakah pertanyaan tersebut masuk akal?
    \item Proses promosi tidak menghasilkan angka apa pun yang bisa
          dijumlahkan. Apa gunanya mencatatnya di gudang data?
    \item Tuliskan satu pertanyaan bisnis yang memerlukan data penjualan
          \emph{dan} persediaan sekaligus.
  \end{enumerate}
\end{enumerate}

\section{Konsep Dasar}

\subsection{Empat Jenis Fact Table}

Jenis fact table ditentukan oleh \emph{perilaku proses}, bukan oleh bentuk
tabel sumbernya \citep{kimball2013}. Pertanyaan penentunya: kapan baris dibuat,
dan apakah baris itu pernah diubah setelah dibuat?

\begin{center}
\small
\begin{tabular}{@{}p{0.20\textwidth}>{\raggedright\arraybackslash}p{0.30\textwidth}>{\raggedright\arraybackslash}p{0.18\textwidth}>{\raggedright\arraybackslash}p{0.24\textwidth}@{}}
\toprule
\textbf{Jenis} & \textbf{Baris dibuat ketika} & \textbf{Diubah?} &
  \textbf{Contoh NusaMart} \\
\midrule
Transaction & Sebuah kejadian terjadi & Tidak pernah &
  Penjualan di kasir \\
Periodic snapshot & Suatu periode berakhir & Tidak pernah &
  Saldo persediaan harian \\
Accumulating snapshot & Sebuah proses dimulai & Berkali-kali &
  Retur barang \\
Factless & Sebuah kejadian atau kondisi berlaku & Tidak pernah &
  Cakupan promosi \\
\bottomrule
\end{tabular}
\end{center}

Ketiga jenis pertama dibedakan oleh hubungannya dengan waktu. Transaction fact
mencatat \emph{titik}; periodic snapshot mencatat \emph{keadaan pada akhir
selang}; accumulating snapshot mengikuti \emph{satu proses yang berjalan} dan
memperbarui barisnya setiap kali proses itu mencapai tonggak berikutnya.

\subsection{Transaction Fact}

Transaction fact mencatat kejadian yang terjadi pada satu titik waktu. Barisnya
dibuat sekali dan tidak pernah diubah --- inilah jenis yang sudah Anda bangun
pada Modul 2 sebagai \perintah{fakta\_penjualan}.

Sifat penting transaction fact: measure-nya \emph{aditif penuh}. Kuantitas dan
subtotal boleh dijumlahkan terhadap dimensi apa pun --- waktu, produk, toko,
atau kombinasi ketiganya --- dan hasilnya selalu bermakna.

\subsection{Periodic Snapshot}

Sebagian proses tidak menghasilkan kejadian, melainkan \emph{keadaan}. Saldo
persediaan adalah contohnya: tidak ada ``transaksi saldo'', yang ada adalah
berapa unit tersisa pada akhir setiap hari.

\begin{konsep}{Measure semi-aditif}
Saldo persediaan boleh dijumlahkan terhadap produk dan toko --- ``berapa total
unit di seluruh toko Lampung pada 3 Januari'' adalah pertanyaan yang sah.
Ia \textbf{tidak} boleh dijumlahkan terhadap waktu: menjumlahkan saldo 1, 2,
dan 3 Januari menghasilkan angka yang tidak mewakili apa pun.

Terhadap waktu, gunakan \perintah{AVG} untuk saldo rata-rata, atau ambil nilai
pada satu titik waktu tertentu dengan window function.
\end{konsep}

Kesalahan ini sangat mudah terjadi karena \perintah{SUM} berjalan tanpa galat
dan menghasilkan angka yang tampak wajar. Saldo satu bulan yang dijumlahkan
akan sekitar tiga puluh kali lipat dari nilai sebenarnya --- cukup besar untuk
salah, tetapi tidak cukup ganjil untuk langsung dicurigai.

Periodic snapshot juga memiliki sifat yang membedakannya dari transaction fact:
\textbf{barisnya padat}, bukan jarang. Produk yang tidak terjual hari ini tidak
menghasilkan baris penjualan, tetapi tetap menghasilkan baris persediaan ---
karena saldonya tetap ada. Konsekuensinya, periodic snapshot sering justru
lebih besar daripada transaction fact.

\subsection{Factless Fact}

Sebagian kejadian penting tidak memiliki measure numerik sama sekali. Bahwa
sebuah produk tercakup sebuah promosi pada suatu hari adalah fakta yang perlu
dicatat, tetapi tidak ada angka yang menyertainya.

\begin{konsep}{Mengapa mencatat yang tidak ada angkanya}
Nilai factless fact terletak pada pertanyaan tentang \emph{ketiadaan}. Data
penjualan hanya memuat produk yang terjual; ia tidak dapat menjawab
``produk mana yang dipromosikan tetapi tidak terjual sama sekali'' --- karena
produk itu tidak muncul di sana.

Factless fact menyediakan daftar \emph{seharusnya-ada}, dan selisihnya terhadap
data penjualan itulah jawabannya.
\end{konsep}

Sebagian factless fact menyertakan satu kolom bernilai tetap 1, disebut
\istilah{counter}, agar \perintah{COUNT} dan \perintah{SUM} berperilaku sama.
Ini pilihan gaya, bukan keharusan; \perintah{COUNT(*)} sudah memadai.

\subsection{Conformed Dimension dan Drill-across}

Dua proses dapat disandingkan hanya bila keduanya memakai dimensi yang
\emph{sama persis} --- tabel yang sama, kunci yang sama, dan makna yang sama.
Dimensi semacam itu disebut \istilah{conformed dimension}, dan inilah yang
Anda tandai pada bus matrix Modul 1.

\begin{peringatan}
Dua tabel fakta \textbf{tidak boleh} di-join langsung pada tingkat detail.
Bila \perintah{fakta\_penjualan} di-join ke
\perintah{fakta\_persediaan\_harian} melalui \perintah{(tanggal, produk,
toko)}, setiap baris penjualan akan berpasangan dengan baris persediaan hari
itu --- dan bila satu produk terjual lima kali dalam sehari, saldo persediaan
hari itu ikut terhitung lima kali.

Galat tidak muncul. Yang muncul hanya angka yang salah.
\end{peringatan}

Cara yang benar disebut \istilah{drill-across}: agregasikan \emph{setiap} fakta
lebih dahulu sampai ke tingkat yang sama, baru sandingkan hasilnya.

\begin{konsep}{Tiga langkah drill-across}
\begin{enumerate}[leftmargin=1.4em]
  \item Tentukan tingkat penyandingan --- kombinasi atribut dimensi bersama,
        misalnya bulan $\times$ kategori $\times$ provinsi.
  \item Agregasikan setiap fakta secara terpisah sampai tingkat itu, memakai
        fungsi agregat yang sesuai sifat measure-nya.
  \item Sandingkan hasilnya dengan \perintah{FULL OUTER JOIN} atas atribut
        dimensi, bukan atas kunci fakta.
\end{enumerate}
\end{konsep}

\perintah{FULL OUTER JOIN} dipakai, bukan \perintah{INNER JOIN}, justru karena
baris yang hanya muncul di satu sisi biasanya adalah temuan yang paling
menarik: produk yang ada stoknya tetapi tidak terjual sama sekali, atau
sebaliknya.

\subsection{Accumulating Snapshot}

Jenis keempat mengikuti satu proses yang melewati beberapa tonggak. Retur
barang di NusaMart melalui empat tahap: pengajuan, persetujuan, penerimaan
barang, dan pengembalian dana.

Yang membedakannya dari tiga jenis lain: \textbf{barisnya diperbarui}. Satu
baris dibuat ketika retur diajukan, dengan sebagian besar kolom tanggal masih
kosong, lalu kolom itu terisi satu per satu seiring proses berjalan. Ini
melanggar kebiasaan ``fact table hanya di-insert'', dan pelanggaran itu
disengaja --- jenis ini memang dirancang untuk mengukur \emph{lama} setiap
tahap.

Perancangan accumulating snapshot untuk proses retur menjadi tugas luar
laboratorium modul ini.

\section{Studi Kasus dan Protokol Praktikum}

Mahasiswa memperluas warehouse NusaMart dari satu fakta penjualan menjadi
beberapa proses yang tetap memakai dimensi terkonformasi.

\begin{center}
\small
\begin{tabular}{@{}p{0.26\textwidth}p{0.20\textwidth}>{\raggedright\arraybackslash}p{0.44\textwidth}@{}}
\toprule
\textbf{Tabel fakta} & \textbf{Jenis} & \textbf{Grain} \\
\midrule
\perintah{fakta\_penjualan} & Transaction &
  Satu produk pada satu baris struk di satu toko pada satu waktu \\
\perintah{fakta\_persediaan\_harian} & Periodic snapshot &
  Satu produk di satu toko pada akhir satu hari \\
\perintah{fakta\_cakupan\_promosi} & Factless &
  Satu produk yang tercakup satu promosi pada satu hari \\
\bottomrule
\end{tabular}
\end{center}

Ketiganya memakai \perintah{dim\_tanggal} dan \perintah{dim\_produk} yang sama;
dua yang pertama juga memakai \perintah{dim\_toko}. Justru dimensi bersama
inilah yang membuat penyandingan pada Bagian D menjadi mungkin.

\begin{catatanasisten}
Fakta cakupan promosi berpotensi besar: jumlah barisnya adalah jumlah produk
per promosi dikalikan lama promosi. Pada dataset kecil, batasi pemuatan pada
satu tahun kalender agar tetap di bawah beberapa ratus ribu baris. Bila
laboratorium terasa lambat, kurangi lagi menjadi satu semester --- pemahaman
konsepnya tidak bergantung pada volume.
\end{catatanasisten}

\section{Alur Praktikum 120 Menit}

\begin{alurwaktu}
0--15 & Demonstrasi pemilihan jenis fact table. \\
15--95 & Kerja terbimbing membangun tiga fakta dan query drill-across. \\
95--110 & Checkpoint grain dan hasil agregasi lintas fakta. \\
110--120 & Penjelasan tugas accumulating snapshot. \\
\end{alurwaktu}

\section{Prosedur Praktikum Terbimbing}

\subsection{Bagian A: Memastikan Grain Setiap Proses}

\langkah{A.1 Menulis tiga pernyataan grain}{10 menit}

Sebelum menulis DDL apa pun, tuliskan grain ketiga fakta pada
\berkas{grain.md}, masing-masing dalam satu kalimat bisnis. Uji setiap kalimat
dengan pertanyaan yang sama seperti Modul 1: dapatkah seorang manajer toko
memahaminya tanpa penjelasan tambahan?

Lalu untuk setiap fakta, tetapkan jenisnya beserta alasan berdasarkan dua
pertanyaan penentu --- kapan baris dibuat, dan apakah baris itu pernah diubah.

\langkah{A.2 Melengkapi junk dimension pada fakta penjualan}{5 menit}

Modul 3 membangun \perintah{dim\_transaksi\_flag}, tetapi
\perintah{fakta\_penjualan} belum menunjuknya. Lengkapi sekarang.

\begin{lstlisting}[style=sqlgaya]
ALTER TABLE fakta_penjualan ADD COLUMN flag_key INTEGER;

UPDATE fakta_penjualan f
SET    flag_key = fl.flag_key
FROM   staging_penjualan s
JOIN   dim_transaksi_flag fl
       ON  fl.metode_bayar = s.metode_bayar
       AND fl.status       = s.status
       AND fl.jenis_pembeli = CASE WHEN s.pelanggan_id IS NULL
                                   THEN 'umum' ELSE 'anggota' END
WHERE  f.transaksi_id = s.transaksi_id
  AND  f.produk_key   = (SELECT dp.produk_key FROM dim_produk dp
                         WHERE dp.produk_id = s.produk_id
                           AND s.waktu_transaksi::DATE >= dp.mulai_berlaku
                           AND s.waktu_transaksi::DATE <  dp.selesai_berlaku);

ALTER TABLE fakta_penjualan ALTER COLUMN flag_key SET NOT NULL;
ALTER TABLE fakta_penjualan ADD CONSTRAINT fk_fakta_flag
  FOREIGN KEY (flag_key) REFERENCES dim_transaksi_flag(flag_key);
\end{lstlisting}

\subsection{Bagian B: Membangun Transaction dan Periodic Snapshot}

\langkah{B.1 Menyalin sumber persediaan ke staging}{5 menit}

Jalankan skrip \berkas{sql/\allowbreak modul-04/\allowbreak
01-staging-persediaan.sql}. Sesudah itu pindahkan isi tabel
\perintah{persediaan\_harian}, \perintah{promosi}, dan
\perintah{promosi\_produk} dari sumber dengan cara yang sama seperti Modul 2.

\langkah{B.2 Membangun periodic snapshot}{15 menit}

\begin{lstlisting}[style=sqlgaya]
-- Grain: satu baris mewakili satu produk di satu toko pada akhir satu hari.
CREATE TABLE fakta_persediaan_harian (
  tanggal_key  INTEGER NOT NULL REFERENCES dim_tanggal(tanggal_key),
  produk_key   INTEGER NOT NULL REFERENCES dim_produk(produk_key),
  toko_key     INTEGER NOT NULL REFERENCES dim_toko(toko_key),
  saldo_awal   NUMERIC(12,2) NOT NULL,   -- semi-aditif
  masuk        NUMERIC(12,2) NOT NULL,   -- aditif
  keluar       NUMERIC(12,2) NOT NULL,   -- aditif
  saldo_akhir  NUMERIC(12,2) NOT NULL,   -- semi-aditif
  PRIMARY KEY (tanggal_key, produk_key, toko_key)
);
\end{lstlisting}

\begin{catatan}
Perhatikan primary key gabungan atas ketiga kunci dimensi. Ini mungkin
dilakukan karena grain periodic snapshot memang menjamin keunikan kombinasi
itu --- satu produk di satu toko hanya punya satu saldo per hari. Pada
transaction fact hal ini tidak berlaku: satu produk dapat terjual beberapa kali
di toko yang sama pada hari yang sama.

Primary key ini sekaligus menjadi penjaga grain: bila pemuatan Anda
menghasilkan dua baris untuk kombinasi yang sama, basis data langsung
menolaknya.
\end{catatan}

Muat dengan pencarian kunci SCD-2 yang sama seperti Modul 3 --- versi dimensi
yang berlaku pada tanggal snapshot, bukan versi kini.

\langkah{B.3 Membuktikan sifat semi-aditif}{10 menit}

Jalankan kedua query berikut untuk satu produk di satu toko selama satu bulan,
lalu bandingkan hasilnya.

\begin{lstlisting}[style=sqlgaya]
-- Keliru: menjumlahkan saldo lintas waktu.
SELECT SUM(saldo_akhir) AS total_yang_menyesatkan
FROM   fakta_persediaan_harian f
JOIN   dim_tanggal dt ON dt.tanggal_key = f.tanggal_key
WHERE  dt.tahun_bulan = 202406 AND f.produk_key = <satu produk>;

-- Benar: rata-rata saldo, dan saldo pada akhir periode.
SELECT ROUND(AVG(f.saldo_akhir), 2) AS rata_saldo,
       (ARRAY_AGG(f.saldo_akhir ORDER BY dt.tanggal DESC))[1] AS saldo_akhir_bulan
FROM   fakta_persediaan_harian f
JOIN   dim_tanggal dt ON dt.tanggal_key = f.tanggal_key
WHERE  dt.tahun_bulan = 202406 AND f.produk_key = <satu produk>;
\end{lstlisting}

Catat kedua angka pada \berkas{catatan-additivity.md} beserta rasionya. Rasio
itu akan mendekati jumlah hari dalam bulan tersebut --- dan itulah besarnya
kesalahan yang ditimbulkan oleh satu \perintah{SUM} yang keliru.

\subsection{Bagian C: Membangun Factless Fact}

\langkah{C.1 Menulis DDL factless fact}{10 menit}

\begin{lstlisting}[style=sqlgaya]
-- Grain: satu baris mewakili satu produk yang tercakup satu promosi
--        pada satu hari.
CREATE TABLE fakta_cakupan_promosi (
  tanggal_key  INTEGER NOT NULL REFERENCES dim_tanggal(tanggal_key),
  produk_key   INTEGER NOT NULL REFERENCES dim_produk(produk_key),
  promosi_key  INTEGER NOT NULL REFERENCES dim_promosi(promosi_key),
  PRIMARY KEY (tanggal_key, produk_key, promosi_key)
);
\end{lstlisting}

Tidak ada satu pun kolom measure. Seluruh isinya adalah kunci dimensi, dan
itulah yang membuatnya disebut \emph{factless}.

\langkah{C.2 Memuat cakupan promosi}{10 menit}

Pemuatan memerlukan pembentangan rentang promosi menjadi satu baris per hari.
Di sini \perintah{generate\_series} kembali dipakai, kali ini untuk membentang
rentang, bukan untuk membangkitkan kalender.

\begin{lstlisting}[style=sqlgaya]
INSERT INTO fakta_cakupan_promosi (tanggal_key, produk_key, promosi_key)
SELECT dt.tanggal_key, dp.produk_key, dpr.promosi_key
FROM   staging_promosi sp
JOIN   staging_promosi_produk spp ON spp.promosi_id = sp.promosi_id
JOIN   dim_promosi dpr            ON dpr.promosi_id = sp.promosi_id
CROSS  JOIN LATERAL generate_series(sp.mulai, sp.selesai, INTERVAL '1 day') g
JOIN   dim_tanggal dt ON dt.tanggal = g::DATE
JOIN   dim_produk  dp ON dp.produk_id = spp.produk_id
                     AND g::DATE >= dp.mulai_berlaku
                     AND g::DATE <  dp.selesai_berlaku
WHERE  EXTRACT(YEAR FROM sp.mulai) = 2024
ON CONFLICT DO NOTHING;
\end{lstlisting}

\begin{catatan}
Klausa \perintah{ON CONFLICT DO NOTHING} menangani produk yang tercakup dua
promosi yang masa berlakunya beririsan. Tanpanya, primary key akan ditolak.
Apakah ini penanganan yang benar bergantung pada aturan bisnis --- bila kedua
promosi memang harus tercatat, grainnya perlu memuat promosi, dan memang itulah
yang dilakukan DDL di atas.
\end{catatan}

\langkah{C.3 Menjawab pertanyaan tentang ketiadaan}{5 menit}

Inilah alasan factless fact dibangun.

\begin{lstlisting}[style=sqlgaya]
-- Produk yang dipromosikan tetapi tidak terjual sama sekali selama promosi.
SELECT dp.nama_produk, dpr.nama_promosi, COUNT(*) AS hari_dipromosikan
FROM   fakta_cakupan_promosi c
JOIN   dim_produk  dp  ON dp.produk_key  = c.produk_key
JOIN   dim_promosi dpr ON dpr.promosi_key = c.promosi_key
WHERE  NOT EXISTS (
         SELECT 1 FROM fakta_penjualan f
         WHERE f.produk_key  = c.produk_key
           AND f.tanggal_key = c.tanggal_key)
GROUP  BY dp.nama_produk, dpr.nama_promosi
ORDER  BY hari_dipromosikan DESC
LIMIT  20;
\end{lstlisting}

Pertanyaan ini mustahil dijawab dari \perintah{fakta\_penjualan} saja: produk
yang tidak terjual tidak memiliki baris di sana.

\subsection{Bagian D: Menulis Query Drill-across}

\langkah{D.1 Membuktikan bahaya join langsung}{10 menit}

Jalankan lebih dahulu cara yang \emph{keliru}, agar Anda melihat akibatnya
sendiri.

\begin{lstlisting}[style=sqlgaya]
-- KELIRU: dua fakta di-join pada tingkat detail.
SELECT SUM(f.kuantitas) AS unit_terjual,
       SUM(p.saldo_akhir) AS saldo_yang_menggelembung
FROM   fakta_penjualan f
JOIN   fakta_persediaan_harian p
       ON  p.tanggal_key = f.tanggal_key
       AND p.produk_key  = f.produk_key
       AND p.toko_key    = f.toko_key;
\end{lstlisting}

Bandingkan \perintah{saldo\_yang\_menggelembung} dengan
\perintah{SUM(saldo\_akhir)} yang dihitung langsung dari
\perintah{fakta\_persediaan\_harian}. Selisihnya adalah akibat penggandaan:
setiap baris persediaan terhitung sebanyak transaksi produk itu pada hari itu.

\langkah{D.2 Menulis drill-across yang benar}{15 menit}

\begin{lstlisting}[style=sqlgaya]
WITH penjualan AS (
  SELECT dt.tahun_bulan, dp.kategori, dtk.provinsi,
         SUM(f.kuantitas) AS unit_terjual,
         SUM(f.subtotal)  AS nilai_penjualan
  FROM   fakta_penjualan f
  JOIN   dim_tanggal dt  ON dt.tanggal_key = f.tanggal_key
  JOIN   dim_produk  dp  ON dp.produk_key  = f.produk_key
  JOIN   dim_toko    dtk ON dtk.toko_key   = f.toko_key
  GROUP  BY 1, 2, 3
),
persediaan AS (
  SELECT dt.tahun_bulan, dp.kategori, dtk.provinsi,
         ROUND(AVG(p.saldo_akhir), 2) AS rata_saldo   -- AVG, bukan SUM
  FROM   fakta_persediaan_harian p
  JOIN   dim_tanggal dt  ON dt.tanggal_key = p.tanggal_key
  JOIN   dim_produk  dp  ON dp.produk_key  = p.produk_key
  JOIN   dim_toko    dtk ON dtk.toko_key   = p.toko_key
  GROUP  BY 1, 2, 3
)
SELECT COALESCE(j.tahun_bulan, s.tahun_bulan) AS tahun_bulan,
       COALESCE(j.kategori,    s.kategori)    AS kategori,
       COALESCE(j.provinsi,    s.provinsi)    AS provinsi,
       j.unit_terjual,
       s.rata_saldo,
       ROUND(j.unit_terjual / NULLIF(s.rata_saldo, 0), 2) AS perputaran
FROM       penjualan  j
FULL OUTER JOIN persediaan s
       ON  s.tahun_bulan = j.tahun_bulan
       AND s.kategori    = j.kategori
       AND s.provinsi    = j.provinsi
ORDER  BY perputaran DESC NULLS LAST;
\end{lstlisting}

Perhatikan tiga hal. Setiap fakta diagregasi \emph{sebelum} disandingkan.
Measure semi-aditif memakai \perintah{AVG}, bukan \perintah{SUM}. Penyandingan
memakai atribut dimensi, bukan kunci fakta.

\begin{catatan}
Baris yang \perintah{unit\_terjual}-nya \perintah{NULL} berarti ada stok tetapi
tidak ada penjualan; baris yang \perintah{rata\_saldo}-nya \perintah{NULL}
berarti ada penjualan tanpa catatan persediaan --- dan yang kedua ini adalah
temuan kualitas data yang layak dilaporkan. Inilah sebabnya
\perintah{FULL OUTER JOIN} dipakai: \perintah{INNER JOIN} akan menyembunyikan
keduanya.
\end{catatan}

\begin{luaran}
\berkas{modul-04-multi-fakta/grain.md}, \berkas{ddl-fakta.sql},
\berkas{load-fakta.sql}, \berkas{catatan-additivity.md}, dan
\berkas{drill-across.sql} beserta \berkas{laporan-drill-across.md}.
\end{luaran}

\begin{checkpoint}
Asisten menjalankan \berkas{sql/modul-04/check.sql} dan memeriksa butir berikut:
\begin{ceklis}
  \item ketiga tabel fakta ada beserta foreign key ke dimensi bersama;
  \item \perintah{fakta\_persediaan\_harian} memiliki primary key gabungan yang
        menjaga grainnya;
  \item \perintah{fakta\_cakupan\_promosi} tidak memuat satu pun kolom measure;
  \item \perintah{fakta\_penjualan} sudah menunjuk \perintah{dim\_transaksi\_flag};
  \item tidak ada baris fakta yang menunjuk versi dimensi di luar masa
        berlakunya;
  \item \berkas{grain.md} memuat tiga pernyataan grain beserta jenis fakta dan
        alasannya;
  \item query drill-across mengagregasi setiap fakta lebih dahulu dan memakai
        \perintah{AVG} untuk measure semi-aditif.
\end{ceklis}
\end{checkpoint}

\section{Latihan Mandiri di Kelas}

\latihan{Mengklasifikasikan proses NusaMart}

Untuk setiap proses berikut, tentukan jenis fact table yang tepat beserta
grainnya. Jawaban tanpa alasan tidak memperoleh nilai penuh.

\begin{enumerate}[leftmargin=1.4em]
  \item Pengiriman barang dari gudang pusat ke toko, yang melewati tahap
        dipesan, dikirim, dan diterima.
  \item Jumlah pelanggan yang mengunjungi toko setiap jam, tercatat oleh
        penghitung di pintu masuk.
  \item Penugasan kasir ke shift kerja.
  \item Saldo poin loyalitas setiap anggota pada akhir bulan.
\end{enumerate}

\latihan{Menakar ukuran periodic snapshot}

Persediaan dicatat untuk setiap produk di setiap toko setiap hari, terlepas
dari ada tidaknya pergerakan.

\begin{enumerate}[leftmargin=1.4em]
  \item Dengan 180.000 produk, 240 toko, dan 365 hari, berapa baris yang
        dihasilkan dalam setahun bila seluruh kombinasi dicatat?
  \item Bandingkan dengan jumlah baris \perintah{fakta\_penjualan} Anda. Mana
        yang lebih besar?
  \item Usulkan satu cara mengurangi ukuran itu tanpa kehilangan kemampuan
        menjawab pertanyaan persediaan, lalu sebutkan apa yang dikorbankan.
\end{enumerate}

\section{Tugas Individu}

\subsection{Pekerjaan Wajib}
Rancang accumulating snapshot untuk proses retur dan jelaskan measure yang
boleh serta tidak boleh disimpan.

Kumpulkan \berkas{accumulating-retur.md} yang memuat:

\begin{enumerate}[leftmargin=1.4em]
  \item pernyataan grain --- perhatikan bahwa satu baris mewakili satu
        \emph{proses}, bukan satu kejadian;
  \item DDL lengkap, memuat satu kunci tanggal untuk \emph{setiap} tonggak,
        seluruhnya menunjuk \perintah{dim\_tanggal} lewat view peran
        masing-masing;
  \item daftar measure lag antar-tonggak beserta cara menghitungnya;
  \item penjelasan measure mana yang \textbf{tidak boleh} disimpan beserta
        alasannya;
  \item penjelasan bagaimana baris diperbarui saat proses berjalan, dan apa
        yang membedakannya dari pemuatan transaction fact.
\end{enumerate}

\begin{catatan}
Butir keempat adalah inti tugas. Pertimbangkan sekurang-kurangnya tiga hal:
rasio dan persentase yang tidak boleh disimpan sebagai measure, status proses
yang berubah dan karenanya bukan measure melainkan atribut dimensi, serta lag
yang belum dapat dihitung karena tonggaknya belum tercapai --- apakah kolomnya
diisi \perintah{NULL}, nol, atau nilai penjaga?
\end{catatan}

\subsection{Pertanyaan Analisis}

\begin{enumerate}[leftmargin=1.4em]
  \item Seorang rekan mengusulkan menggabungkan penjualan dan persediaan
        menjadi satu tabel fakta agar tidak perlu drill-across. Jelaskan
        mengapa usulan ini keliru, dengan menunjuk pada grain kedua proses.
  \item Pada query drill-across Anda, sebagian baris memiliki
        \perintah{rata\_saldo} tetapi \perintah{unit\_terjual}-nya
        \perintah{NULL}. Apa arti bisnisnya, dan mengapa baris ini justru lebih
        menarik daripada baris yang lengkap?
  \item Factless fact cakupan promosi Anda memuat ratusan ribu baris tanpa satu
        pun measure. Seorang rekan menyebutnya pemborosan ruang. Bantah dengan
        satu pertanyaan bisnis yang hanya dapat dijawab olehnya.
\end{enumerate}

\section{Format Pengumpulan dan Checklist}

Kumpulkan DDL tiga fakta, pernyataan grain, query drill-across, dan rancangan
accumulating snapshot.

Seluruh berkas diletakkan pada \berkas{modul-04-multi-fakta/}.

\begin{ceklis}
  \item \berkas{pre-lab.md}
  \item \berkas{grain.md} --- tiga grain, jenis fakta, dan alasannya
  \item \berkas{ddl-fakta.sql}
  \item \berkas{load-fakta.sql}
  \item \berkas{catatan-additivity.md} --- perbandingan \perintah{SUM} dan
        \perintah{AVG} beserta rasionya
  \item \berkas{drill-across.sql} dan \berkas{laporan-drill-across.md}
  \item \berkas{accumulating-retur.md}
  \item \berkas{jawaban-analisis.md}
\end{ceklis}

\section{Troubleshooting}

\begin{table}[htbp]
\centering
\small
\caption{Masalah yang lazim muncul pada Modul 4 dan penanganannya}
\label{tab:m4-troubleshooting}
\begin{tabular}{@{}p{0.30\textwidth}>{\raggedright\arraybackslash}p{0.62\textwidth}@{}}
\toprule
\textbf{Gejala} & \textbf{Penanganan} \\
\midrule
Angka persediaan jauh lebih besar daripada yang wajar &
  Measure semi-aditif dijumlahkan lintas waktu. Ganti \perintah{SUM} dengan
  \perintah{AVG}, atau ambil nilai pada satu titik waktu. \\
Total measure berubah setelah menambah join ke fakta kedua &
  Dua tabel fakta di-join pada tingkat detail. Agregasikan masing-masing lebih
  dahulu, lalu sandingkan hasilnya. \\
\perintah{duplicate key value} saat memuat periodic snapshot &
  Sumber memuat lebih dari satu baris untuk kombinasi hari, produk, dan toko
  yang sama --- grain sumber tidak sesuai dugaan. Periksa dengan
  \perintah{GROUP BY ... HAVING COUNT(*) > 1}. \\
Factless fact membengkak jauh melebihi perkiraan &
  Rentang promosi dibentang tanpa batas tahun, atau produk tercakup beberapa
  promosi sekaligus. Batasi tahun pemuatan dan periksa
  \perintah{staging\_promosi\_produk}. \\
Drill-across kehilangan banyak baris &
  \perintah{INNER JOIN} dipakai, sehingga kombinasi yang hanya ada di satu
  fakta hilang. Ganti dengan \perintah{FULL OUTER JOIN} dan
  \perintah{COALESCE} pada kolom kuncinya. \\
Kolom kunci drill-across bernilai \perintah{NULL} &
  \perintah{COALESCE} belum dipasang pada atribut penyandingan. Kunci harus
  diambil dari sisi mana pun yang tersedia. \\
Pembagian menghasilkan galat \emph{division by zero} &
  Gunakan \perintah{NULLIF(penyebut, 0)} agar hasilnya \perintah{NULL},
  bukan galat. \\
\bottomrule
\end{tabular}
\end{table}

\section{Ringkasan}

Modul ini menambah dua proses ke gudang data dan, bersamaan dengan itu, satu
aturan yang berlaku selamanya: \textbf{dua tabel fakta tidak pernah di-join
pada tingkat detail}. Yang disandingkan adalah hasil agregat, bukan baris.

Tiga gagasan perlu dibawa ke Modul 5. \emph{Pertama}, jenis fact table
ditentukan oleh perilaku proses --- kapan barisnya dibuat dan apakah pernah
diubah --- bukan oleh bentuk tabel sumbernya. \emph{Kedua}, sifat keteraditifan
measure adalah bagian dari rancangan, bukan urusan query: measure semi-aditif
yang dijumlahkan lintas waktu menghasilkan angka yang salah tanpa galat apa
pun. \emph{Ketiga}, dimensi terkonformasi bukan kerapian melainkan syarat ---
tanpa dimensi yang benar-benar sama, dua proses tidak dapat disandingkan sama
sekali.

Gudang data Anda kini memuat tiga tabel fakta dan enam dimensi, seluruhnya
masih pada schema \perintah{public} dengan tipe data yang dipilih seadanya.
Modul 5 menata ulang keseluruhannya menjadi objek fisik yang dapat
dipertanggungjawabkan byte-nya --- termasuk menghitung lebar baris ketiga fakta
yang baru saja Anda bangun.
